\documentclass[10pt,a4paper]{article}
	
	\usepackage{amsmath,amssymb,amsthm,mathtools,amsfonts}
	\usepackage{breqn}
	\usepackage[utf8]{inputenc}
	\usepackage[english]{babel}
	\usepackage[T1]{fontenc}
	
	%%%%%%%%%%%%%%%%%%%%%%%%%%%%%%%%%%%%%%%%%%%%%%%%%%%%%%%%%%%%%%%%%%%%%%%%%%%%%% Fonts
	
	\usepackage{fontspec}
	\setmainfont[Numbers=OldStyle]{EB Garamond}
	\newfontfamily\plantinc[Numbers=OldStyle]{Plantin MT Pro Bold Condensed}
	\newfontfamily\plantinsb{Plantin MT Pro Semibold}
	\newfontfamily\garamond[Numbers=OldStyle]{EB Garamond}
	\newfontfamily\plantin[Numbers=OldStyle]{Plantin MT Pro}
	\newfontfamily\wal{Walbaum Com}
	\newfontfamily\klav{Klavika}
	\newfontfamily\klavl{Klavika light}
	\newfontfamily\quotefont[Ligatures=TeX]{Linux Libertine O}
	
	%%%%%%%%%%%%%%%%%%%%%%%%%%%%%%%%%%%%%%%%%%%%%%%%%%%%%%%%%%%%%%%%%%%%%%%%%%%%%% Colors
	
	\usepackage[svgnames]{xcolor}
	
	\definecolor{hgblue}{RGB}{10, 40, 80}
	\definecolor{hgred}{RGB}{140, 10, 10}
	\definecolor{hggrey}{RGB}{215, 215, 210}
	\definecolor{hggreen}{RGB}{145, 205, 205}
	
	%%%%%%%%%%%%%%%%%%%%%%%%%%%%%%%%%%%%%%%%%%%%%%%%%%%%%%%%%%%%%%%%%%%%%%%%%%%%%% Other Graphics
	
	\usepackage{graphicx}
	\graphicspath{{./img/}}
	\usepackage{tikz}
	\usepackage{placeins}
	\usepackage[modulo]{lineno}
	%\usepackage{geometry}
	
	\usepackage[pdfborder={0 0 0}]{hyperref} %% Ingen firkant rundt om referencer

	\usepackage{multicol}
	\usepackage{cellspace}
		\setlength\cellspacetoplimit{100pt}
		\setlength\cellspacebottomlimit{100pt}
	\usepackage{tabularx}
		\newcolumntype{b}{>{\hsize=\hsize}>{\centering\arraybackslash}X}
	\usepackage{siunitx}
	\usepackage{cellspace}
		\setlength\cellspacetoplimit{4pt}
		\setlength\cellspacebottomlimit{4pt}
	
	\usepackage{enumitem}% better controls with enumerating
	\usepackage{wrapfig}
		%\setlength{\wrapoverhang}{\marginparwidth}
		%\addtolength{\wrapoverhang}{\marginparsep}
		\setlength{\columnsep}{0pt}
		\setlength{\intextsep}{-10pt}
	\usepackage{caption}
	%\usepackage{dirtytalk}
	%\usepackage{tasks}
	
	

	%\setlength{\parskip}{0.5\baselineskip}
	%\setlength{\parindent}{0cm}
	
	\setlist[enumerate]{label=\alph*),left=20pt}
	%\setlist[enumerate]{leftmargin=\parindent}
	%\setlist[enumerate]{nosep}
	
	\usepackage{lipsum}
	\usepackage{comment}
	%\usepackage{ulem}
	
	%\reversemarginpar
	
	%%%%%%%%%%%%%%%%%%%%%%%%%%%%%%%%%%%%%%%%%%%%%%%%%%%%%%%%%%%%%%%%%%%%%%%%%%%%%% New commands
	
	\newcommand\svar[1]{\newline\textcolor{red}{\textbf{Svar}\\#1}}
	
	\usepackage[h]{esvect}
	\newcommand{\twvect}[2]{\ensuremath{\begin{pmatrix}#1\\#2\end{pmatrix}}}
	
	\newcommand*\quotesize{20}
	\newcommand*{\openquote}
		{\tikz[remember picture,overlay,xshift=-3ex,yshift=-0ex]
			\node (OQ) {\quotefont\fontsize{\quotesize}{\quotesize}\selectfont``};\kern0pt}

	\newcommand*{\closequote}[1]
		{\tikz[remember picture,overlay,xshift=2ex,yshift={#1}]
			\node (CQ) {\quotefont\fontsize{\quotesize}{\quotesize}\selectfont''};}
			
	\newcommand{\qm}[1]{‘‘#1”}
	\newcommand{\iquote}[1]{\begin{quote}\itshape \openquote#1\hfill\closequote{0ex} \end{quote}}

	\newcommand{\source}[1]{\footnote{Source: \href{#1}{#1}}}
	
	\newcommand{\glos}[3]{\dotuline{#1}\marginpar{\scriptsize\textit{#3} #2}}
	
	\newcommand{\subtitle}[1]{\vspace{10pt}\newline\normalsize\plantin #1}
	
	%%%%%%%%%%%%%%%%%%%%%%%%%%%%%%%%%%%%%%%%%%%%%%%%%%%%%%%%%%%%%%%%%%%%%%%%%%%%%% Sections
	
	\usepackage{titlesec}
	\titleformat{\subsection}{\color{hgred}\plantinc\large}{}{}{}
	
	\usepackage{titling}
	\setlength{\droptitle}{-12ex}
	\pretitle{\begin{flushleft}\plantinc\huge\color{hgblue}}
	\posttitle{\par\end{flushleft}}
	\preauthor{\begin{flushleft}\Large}
	\postauthor{\end{flushleft}}
	\predate{\begin{flushleft}}
	\postdate{\end{flushleft}}
	
	\setcounter{secnumdepth}{0}
	
	%%%%%%%%%%%%%%%%%%%%%%%%%%%%%%%%%%%%%%%%%%%%%%%%%%%%%%%%%%%%%%%%%%%%%%%%%%%%%% Header & Footer
	
	\usepackage{bigfoot}
	\DeclareNewFootnote{a}
		\renewcommand{\thefootnotea}{\fnsymbol{footnotea}}
			\newcommand{\footnotes}[1]{\footnotea[2]{#1}} % here you can specify what symbol you want in footnotes
			\newcommand{\footnoted}[1]{\footnotea[3]{#1}} % for a second footnote on a page
			\MakePerPage{footnotes}
	
			%%%%%%%%%%%%%%%%%%%%%%%%%%%%%%%%%%%%%%%%%%%%%%%%%%%%%%%%%%%%%%%%%%%%%%%%%%%%%% Document

\begin{document}
%\pagecolor{FloralWhite}
\title{Girl}
%\date{\vspace{-3em}} % I tilfældet ingen dato
\date{1978}
\author{Jamaica Kincaid}
\maketitle
    
\thispagestyle{plain}
\setcounter{page}{1}
\linenumbers
\noindent Wash the white clothes on Monday and put them on the stone heap; wash the color clothes on Tuesday and put them on the clothesline to dry; don’t walk bare-head in the hot sun; cook pumpkin fritters in very hot sweet oil; soak your little cloths right after you take them off; when buying cotton to make yourself a nice blouse, be sure that it doesn’t have gum in it, because that way it won’t hold up well after a wash; soak salt fish overnight before you cook it; is it true that you sing benna in Sunday school?; always eat your food in such a way that it won’t turn someone else’s stomach; on Sundays try to walk like a lady and not like the slut you are so bent on becoming; don’t sing benna in Sunday school; you mustn’t speak to wharf-rat boys, not even to give directions; don’t eat fruits on the street—flies will follow you; \emph{but I don’t sing benna on Sundays at all and never in Sunday school}; this is how to sew on a button; this is how to make a buttonhole for the button you have just sewed on; this is how to hem a dress when you see the hem coming down and so to prevent yourself from looking like the slut I know you are so bent on becoming; this is how you iron your father’s khaki shirt so that it doesn’t have a crease; this is how you iron your father’s khaki pants so that they don’t have a crease; this is how you grow okra—far from the house, because okra tree harbors red ants; when you are growing dasheen, make sure it gets plenty of water or else it makes your throat itch when you are eating it; this is how you sweep a corner; this is how you sweep a whole house; this is how you sweep a yard; this is how you smile to someone you don’t like too much; this is how you smile to someone you don’t like at all; this is how you smile to someone you like completely; this is how you set a table for tea; this is how you set a table for dinner; this is how you set a table for dinner with an important guest; this is how you set a table for lunch; this is how you set a table for breakfast; this is how to behave in the presence of men who don’t know you very well, and this way they won’t recognize immediately the slut I have warned you against becoming; be sure to wash every day, even if it is with your own spit; don’t squat down to play marbles—you are not a boy, you know; don’t pick people’s flowers—you might catch something; don’t throw stones at blackbirds, because it might not be a blackbird at all; this is how to make a bread pudding; this is how to make doukona; this is how to make pepper pot; this is how to make a good medicine for a cold; this is how to make a good medicine to throw away a child before it even becomes a child; this is how to catch a fish; this is how to throw back a fish you don’t like, and that way something bad won’t fall on you; this is how to bully a man; this is how a man bullies you; this is how to love a man, and if this doesn’t work there are other ways, and if they don’t work don’t feel too bad about giving up; this is how to spit up in the air if you feel like it, and this is how to move quick so that it doesn’t fall on you; this is how to make ends meet; always squeeze bread to make sure it’s fresh; \emph{but what if the baker won’t let me feel the bread?}; you mean to say that after all you are really going to be the kind of woman who the baker won’t let near the bread?








































\end{document}



The villagers kept their distance, leaving a space between themselves and the stool, and when Mr. Summers said, “Some of you fellows want to give me a hand?,” there was a hesitation before two men, Mr. Martin and his oldest son, Baxter, came forward to hold the box steady on the stool while Mr. Summers stirred up the papers inside it.

The original paraphernalia for the lottery had been lost long ago, and the black box now resting on the stool had been put into use even before Old Man Warner, the oldest man in town, was born. Mr. Summers spoke frequently to the villagers about making a new box, but no one liked to upset even as much tradition as was represented by the black box. There was a story that the present box had been made with some pieces of the box that had preceded it, the one that had been constructed when the first people settled down to make a village here. Every year, after the lottery, Mr. Summers began talking again about a new box, but every year the subject was allowed to fade off without anything’s being done. The black box grew shabbier each year; by now it was no longer completely black but splintered badly along one side to show the original wood color, and in some places faded or stained.

Mr. Martin and his oldest son, Baxter, held the black box securely on the stool until Mr. Summers had stirred the papers thoroughly with his hand. Because so much of the ritual had been forgotten or discarded, Mr. Summers had been successful in having slips of paper substituted for the chips of wood that had been used for generations. Chips of wood, Mr. Summers had argued, had been all very well when the village was tiny, but now that the population was more than three hundred and likely to keep on growing, it was necessary to use something that would fit more easily into the black box. The night before the lottery, Mr. Summers and Mr. Graves made up the slips of paper and put them into the box, and it was then taken to the safe of Mr. Summers’ coal company and locked up until Mr. Summers was ready to take it to the square next morning. The rest of the year, the box was put away, sometimes one place, sometimes another; it had spent one year in Mr. Graves’ barn and another year underfoot in the post office, and sometimes it was set on a shelf in the Martin grocery and left there.


“Clark. . . . Delacroix.”

“There goes my old man,” Mrs. Delacroix said. She held her breath while her husband went forward.

“Dunbar,” Mr. Summers said, and Mrs. Dunbar went steadily to the box while one of the women said, “Go on, Janey,” and another said, “There she goes.”

“We’re next,” Mrs. Graves said. She watched while Mr. Graves came around from the side of the box, greeted Mr. Summers gravely, and selected a slip of paper from the box. By now, all through the crowd there were men holding the small folded papers in their large hands, turning them over and over nervously. Mrs. Dunbar and her two sons stood together, Mrs. Dunbar holding the slip of paper.

“Harburt. . . . Hutchinson.”

“Get up there, Bill,” Mrs. Hutchinson said, and the people near her laughed.

“Jones.”
